% Part: lambda-calculus
% Chapter: lambda-definability
% Section: lambda-definable-recursive

\documentclass[../../../include/open-logic-section]{subfiles}

\begin{document}

\olfileid{lam}{dfl}{ldr}
\olsection{\usetoken{S}{lambda definable} 函数是递归的}

不仅所有部分递归函数都是 $\lambd$-可定义的，其逆亦真。即，所有 $\lambd$-可定义的函数都是部分递归的。

\begin{thm}
  \ollabel{thm:lambda-computable} 若一个部分函数 $f$ 是 $\lambd$-可定义的，则它是部分递归的。
\end{thm}

\begin{proof}
我们仅概述证明。首先，对 $\lambd$-项进行算术化，即系统地为其赋予哥德尔数，采用通常的素数幂序列编码。然后，定义一个部分递归函数 $\fn{normalize}(t)$，它以某个 $\lambd$-项的哥德尔数~$t$ 为自变元，若该 $\lambd$-项有范式，则返回该范式的哥德尔数，否则无定义。接着，定义两个部分递归函数 $\fn{toChurch}$ 和 $\fn{fromChurch}$，分别将自然数映射为对应丘奇数的哥德尔数，以及将对应丘奇数的哥德尔数映射回自然数。

利用这些递归函数，我们可以将函数~$f$ 定义为一个部分递归函数。存在一个 $\lambd$-项~$F$ $\lambd$-定义了~$f$。为计算 $f(n_1, \dots, n_k)$，首先使用 $\fn{toChurch}(n_i)$ 得到对应丘奇数的哥德尔数，将它们追加到 $\Gn{F}$ 上，得到项 $F \num{n_1}\dots\num{n_k}$ 的哥德尔数。然后，对该哥德尔数应用 $\fn{normalize}$。若 $f(n_1, \dots, n_k)$ 有定义，则 $F \num{n_1}\dots\num{n_k}$ 有范式（该范式必为一个丘奇数），否则它没有范式（因而 \[\fn{normalize}(\Gn{F\num{n_1}\dots\num{n_k}})\] 无定义）。最后，对已化为范式的项的哥德尔数使用 $\fn{fromChurch}$。
\end{proof}

\end{document}